\documentclass[12pt]{article}

\usepackage{geometry}
\usepackage{titling}
\usepackage{amsmath,amscd}
\usepackage{fontspec}
\usepackage{unicode-math}
\usepackage{pifont}
\usepackage{xcolor}
\usepackage[dvipsnames]{xcolor}
\usepackage{listings}
\usepackage{minted}
\usepackage{tikz}
\usepackage{nicematrix}

\usetikzlibrary{fit}

\geometry{a4paper, top=2.5cm, bottom=2.5cm, left=2.2cm, right=2.2cm}
\setlength{\droptitle}{-3cm}

\setmonofont{DejaVu Sans Mono}
\setmainfont{Libertinus Serif}
\setsansfont{Libertinus Sans}
\setmathfont{STIX Two Math}

\newcommand{\vsp}{\vspace*{1em}}
\newcommand{\snl}{\vspace*{2pt}}

\title{Finite-Dimensional Vector Spaces}
\author{Devi Prasad M}
\date{August 2026}
\begin{document}

\maketitle

\section{\emph{Auto-regressive model}}.

Suppose that $z_1,z_2,\ldots$ is a time series, with the number $z_t$ giving
the value in period or time $t$. For example $z_t$ could be the gross sales
at a particular store on day $t$. An \emph{auto-regressive} (AR) model is used to
predict $z_{t+1}$ from the previous
$M$ values, $z_t,z_{t-1},\ldots,z_{t-M+1}$:
\[
    z_{t+1} = (z_t,z_{t-1},\ldots z_{t-M+1})^{\textrm{T}} \beta, \ \, t= M,M + 1,\ldots
\]

Here $z_{t+1}$ denotes the AR model's prediction of $z_{t+1}$, $M$ is the
memory length of the AR model, and the $M$-vector $\beta$ is the AR model
coefficient vector. For this problem we will assume that the time period
is daily, and $M = 10$. Thus, the AR model predicts tomorrow's value,
given the values over the last 10 days.

For each of the following cases, give a short interpretation or description of the AR model
in English, without referring to mathematical concepts like vectors, inner product, and
so on. You can use words like `yesterday' or `today'.

\begin{itemize}
    \item [(a)] $\beta \approx e_1$.
    \item [(b)] $\beta \approx 2e_1 - e_2$.
    \item [(c)] $\beta \approx e_6$.
    \item [(d)] $\beta \approx 0.5 e_1 + 0.5 e_2$.
\end{itemize}

\subsection{Answer}
{\noindent{\textbf(a)} \ \ $\beta \approx e_1$}\snl

Tomorrow's value will be roughly the same as today's value.
The model ignores everything that happened earlier in the past 10 days
and simply predicts that whatever sales or numbers you had today will
repeat tomorrow.

\vsp

{\noindent{\textbf{(b)}} \ \ $\beta \approx 2e_1 - e_2$}\snl

Tomorrow's value continues today's trend (the direction and amount of
change from yesterday to today).

If sales grew by $100$ between yesterday and today, the model predicts
they will grow by another $100$ tomorrow. Expressed as
$e_1 + (e_1 - e_2)$ it takes today's value and adds today's change.

\vsp

{\noindent{\textbf{(c)}} \ \ $\beta \approx e_6$}\snl

Tomorrow's value will match the value from 6 days ago.
The model looks back almost a full week and assumes tomorrow will look
identical to whatever happened on that same day earlier in the cycle.

\vsp

{\noindent{\textbf{(d)}} \ \ $\beta \approx 0.5e_1 + 0.5e_2$}\snl

Tomorrow's value will be the simple average of today's and yesterday's values.
The model smooths out day-to-day noise by taking a two-day moving average and
using that to predict what comes next.

\end{document}